% !TeX encoding = UTF-8
\documentclass[variant=apuntes,cover=institutional,mode=screen]{ua-engineering-v6-article}
% !TeX encoding = UTF-8
\UASetUniversity{Universidad Austral}
\UASetFaculty{Facultad de Ingeniería}
\UASetDepartment{Departamento de Computación}
\UASetCampus{Pilar}
\UASetCareer{Ingeniería Informática}
\UASetCommission{Comisión A}
\UASetProfessor{Equipo docente}
\UASetSemester{Segundo cuatrimestre 2026}
\UASetCourse{Sistemas Distribuidos}
\UASetDate{2026}

\UASetSemester{Segundo cuatrimestre 2026}
\UASetDate{2026}

\UASetClassNumber{09}
\UASetClassTitle{Arquitecturas distribuidas reales}
\title{Clase 09 --- Monolito modular, microservicios, eventos, outbox y sagas}
\author{\UAProfessor}
\date{\UADate}

\begin{document}
\maketitle

\section{Propósito de la clase}

Esta clase cambia el nivel de análisis. Las clases anteriores estudiaron mecanismos de almacenamiento, logs y consistencia. Ahora debemos decidir cómo agrupar responsabilidades, datos y equipos en una arquitectura completa.

El objetivo no es defender una moda arquitectónica. Es aprender a responder:
\begin{itemize}
\item qué frontera de dominio justifica una frontera de despliegue;
\item qué invariantes conviene mantener locales;
\item qué interacciones requieren respuesta inmediata;
\item qué procesos pueden evolucionar mediante eventos;
\item cómo se conserva integridad de negocio cuando no existe una única transacción;
\item si la organización puede operar lo que diseña.
\end{itemize}

\begin{uakeybox}
Una arquitectura distribuida no mejora por contener más procesos. Mejora cuando sus fronteras reducen acoplamiento relevante sin ocultar los nuevos costos de red, datos y operación.
\end{uakeybox}



\section{Cómo estudiar este capítulo}

Este capítulo está organizado en tres pasadas. La primera construye una intuición con un caso real; la segunda formaliza el mecanismo y sus supuestos; la tercera somete la solución a una falla o a una medición. Para estudiar de manera eficaz:

\begin{enumerate}
\item explique primero el problema sin mencionar una tecnología;
\item dibuje la ejecución normal y mueva el punto de falla;
\item identifique la garantía exacta y su frontera;
\item anote qué costo se desplaza hacia latencia, almacenamiento, reparación u operación;
\item cierre proponiendo una métrica o un experimento que podría refutar su diseño.
\end{enumerate}

\begin{uainfo}
Las analogías del capítulo son puertas de entrada, no definiciones finales. Una respuesta completa debe poder abandonar la analogía y utilizar el vocabulario técnico correspondiente.
\end{uainfo}

\section{Mapa intuitivo antes del formalismo}

La arquitectura se entiende mejor si primero se piensa en responsabilidades y no en productos:

\begin{center}
\begin{tabular}{p{0.18\textwidth}p{0.34\textwidth}p{0.36\textwidth}}
\toprule
Concepto & Intuición & Pregunta técnica \\
\midrule
Monolito modular & edificio único con habitaciones y puertas & ¿las paredes internas se respetan? \\
Microservicio & edificio autónomo conectado por caminos & ¿qué autonomía compra la separación? \\
Bounded context & idioma local de una parte del negocio & ¿qué significa aquí ``pedido''? \\
Ownership & escritura registral del dato & ¿quién puede modificar el invariante? \\
RPC & llamada telefónica que espera respuesta & ¿qué ocurre si tarda o no responde? \\
Evento & aviso publicado en un tablero durable & ¿quién lo consume y con qué retraso? \\
Outbox & copia del aviso en el mismo libro que el cambio & ¿cómo evitar estado sin evento? \\
Saga & itinerario con acciones compensatorias & ¿cómo continuar o reparar sin rollback global? \\
\bottomrule
\end{tabular}
\end{center}

\begin{uainfo}
Una arquitectura no es un dibujo de cajas. Es una distribución de autoridad, datos, fallas, comunicación y responsabilidad operacional.
\end{uainfo}

\section{Caso conductor: Asteria Online}

Asteria es un mundo virtual medieval. Cuenta con cuentas, zonas de mundo, combate, inventario, mercado, pagos, ranking y chat. Las presiones son distintas:

\begin{itemize}
\item combate necesita baja latencia y autoridad local;
\item inventario debe preservar propiedad única de objetos;
\item pagos requieren auditabilidad e idempotencia;
\item ranking tolera atraso;
\item chat puede degradarse sin detener la partida;
\item zonas populares necesitan escalar de forma distinta.
\end{itemize}

La primera pregunta no es cuántos servicios crear, sino qué capacidades poseen invariantes, ritmo de cambio, carga y equipo suficientemente distintos para justificar una frontera de despliegue.

\begin{center}
\begin{tikzpicture}[node distance=6mm,>=Latex,every node/.style={font=\small,align=center}]
\node[draw=UAIngOrange,rounded corners,thick,fill=UAIngPaper,text width=2cm] (acc) {Cuentas};
\node[draw=UAIngNavy,rounded corners,thick,right=of acc,text width=2cm] (world) {Mundo};
\node[draw=UAIngOrange,rounded corners,thick,right=of world,text width=2cm] (inv) {Inventario};
\node[draw=UAIngNavy,rounded corners,thick,right=of inv,text width=2cm] (pay) {Pagos};
\node[draw=UAIngOrange,rounded corners,thick,right=of pay,text width=2cm] (rank) {Ranking};
\draw[->,very thick] (acc)--(world);
\draw[->,very thick,UAIngOrange] (world)--(inv);
\draw[->,very thick] (inv)--(pay);
\draw[->,very thick,UAIngOrange] (pay)--(rank);
\end{tikzpicture}
\end{center}

\section{Caso conductor: plataforma global de cursos}

La plataforma ofrece catálogo, inscripción, pagos, reproducción, progreso, recomendaciones y certificados. Existen cinco equipos con ritmos de cambio distintos. Algunas propiedades:
\begin{itemize}
\item inscripción, cupo, pago y certificado tienen invariantes fuertes;
\item recomendaciones y analytics aceptan atraso;
\item video tiene una carga y SLO diferentes;
\item el dominio de aprendizaje todavía cambia;
\item el negocio quiere despliegues más independientes.
\end{itemize}

Una mala respuesta sería “microservicios para escalar”. Una respuesta madura analiza qué capacidades extraer, qué mantener juntas y qué mecanismos sostienen la separación.

\section{Arquitectura como sistema de decisiones}

Una arquitectura define:
\begin{itemize}
\item responsabilidades y límites;
\item ownership de datos;
\item contratos sincrónicos y asincrónicos;
\item consistencia e invariantes;
\item despliegue y escalado;
\item modos de falla y degradación;
\item observabilidad y operación;
\item evolución y migración.
\end{itemize}

El diagrama es evidencia visual de esas decisiones, no la arquitectura en sí.

\section{Monolito modular}

\subsection{Definición}

Una aplicación se despliega como unidad, pero internamente separa módulos con APIs, dependencias y ownership claros. El mismo proceso no implica modelo global desordenado.

\subsection{Ventajas}

\begin{itemize}
\item llamadas locales simples y rápidas;
\item transacciones sobre una base común;
\item debugging y testing end-to-end más accesibles;
\item despliegue y rollback unificados;
\item menor superficie de red y operación;
\item facilidad para refactorizar fronteras mientras el dominio se aprende.
\end{itemize}

\subsection{Riesgos}

\begin{itemize}
\item dependencias internas sin control;
\item cambios que atraviesan módulos;
\item acoplamiento por modelo de datos compartido;
\item despliegues cada vez más grandes;
\item dificultad de ownership si todos modifican todo.
\end{itemize}

El problema no es la unidad de despliegue, sino perder modularidad.

\section{Monolith First y prima de microservicios}

Martin Fowler y otros autores han señalado que comenzar con un monolito modular puede ser razonable cuando las fronteras del dominio aún no están claras. Extraer demasiado pronto congela hipótesis equivocadas en contratos de red y pipelines operativos.

La “prima de microservicios” incluye:
\begin{itemize}
\item timeouts, retries e idempotencia;
\item compatibilidad de contratos;
\item datos distribuidos y sagas;
\item tracing y correlación;
\item service discovery, secretos y configuración;
\item despliegue, rollback y capacidad de plataforma;
\item incidentes que atraviesan equipos.
\end{itemize}

La separación es racional cuando sus beneficios superan esa prima.

\section{Microservicios}

\subsection{Definición operacional}

Un microservicio útil posee una capacidad de negocio coherente, datos y comportamiento bajo un ownership claro, y puede evolucionar con una independencia razonable.

No debe definirse por cantidad de líneas, tamaño de equipo o uso de HTTP.

\subsection{Señales que justifican separación}

\begin{itemize}
\item carga y escalado muy diferentes;
\item ciclos de cambio independientes;
\item equipos con ownership estable;
\item necesidad de aislar fallas;
\item requerimientos de seguridad o cumplimiento distintos;
\item tecnología especializada con beneficio real;
\item frontera de dominio relativamente madura.
\end{itemize}

\subsection{Caso real: Netflix}

Netflix ha documentado plataformas donde componentes separados requieren escalado y despliegue independientes. El ejemplo es valioso si se entiende la causa: la escala y organización justifican la separación. Copiar la forma sin la presión original produce complejidad sin beneficio.

\section{Monolito distribuido}

Un sistema tiene muchos servicios, pero conserva acoplamiento de monolito:
\begin{itemize}
\item despliegues coordinados siempre;
\item base compartida;
\item cadenas sincrónicas largas;
\item modelos y librerías compartidos rígidamente;
\item cambios que requieren múltiples equipos;
\item fallas locales que detienen el flujo completo.
\end{itemize}

Se paga la latencia, observabilidad y operación distribuida sin obtener autonomía.

\section{Bounded contexts}

Un bounded context define una frontera donde términos y reglas tienen significado coherente. La palabra “curso” puede representar catálogo, experiencia de aprendizaje o producto facturable.

Intentar un modelo canónico universal aumenta acoplamiento. Traducir entre contextos mediante contratos hace explícitas las diferencias y reduce el riesgo de que una semántica local invada a otra.

\section{Ownership de datos}

Database per service no significa necesariamente un servidor físico por servicio. Significa que un servicio es el único que modifica directamente su modelo. Otros acceden mediante contratos.

Beneficios:
\begin{itemize}
\item autonomía de esquema;
\item invariantes locales claras;
\item menor acoplamiento accidental;
\item accountability operacional.
\end{itemize}

Costos:
\begin{itemize}
\item joins globales requieren proyecciones o composición;
\item transacciones atraviesan servicios;
\item consistencia se vuelve explícita;
\item reporting necesita pipelines y modelos derivados.
\end{itemize}

\section{Conway y organización}

La arquitectura tiende a reflejar las rutas de comunicación. Si dos equipos deben coordinar cada cambio, una frontera técnica no produce autonomía. Un servicio sin equipo dueño acumula incidentes y deuda.

La pregunta no es solo “¿podemos desplegarlo aparte?”, sino “¿existe un equipo con contexto, on-call y autoridad para operarlo?”.

\section{RPC y eventos}

\subsection{RPC}

Ventajas:
\begin{itemize}
\item control de flujo explícito;
\item respuesta inmediata;
\item errores visibles en el request;
\item adecuado para queries y comandos críticos sincrónicos.
\end{itemize}

Costos:
\begin{itemize}
\item acoplamiento temporal;
\item latencia agregada;
\item cascadas de timeout/retry;
\item disponibilidad compuesta;
\item necesidad de versionado y compatibilidad.
\end{itemize}

\subsection{Eventos}

Ventajas:
\begin{itemize}
\item desacoplamiento temporal;
\item múltiples consumidores;
\item replay y auditoría;
\item extensión sin modificar productor.
\end{itemize}

Costos:
\begin{itemize}
\item consistencia eventual;
\item causalidad y orden parcial;
\item duplicados;
\item trazabilidad más difícil;
\item contratos durables;
\item estados intermedios visibles.
\end{itemize}

No se debe elegir eventos por moda. Una consulta de disponibilidad que necesita respuesta inmediata puede ser mejor por RPC; una notificación o proyección puede ser evento.

\section{Comandos, eventos y consultas}

\begin{itemize}
\item Un comando expresa intención a un owner: \texttt{ReserveSeat}.
\item Un evento expresa un hecho: \texttt{SeatReserved}.
\item Una consulta solicita información sin mutar: \texttt{GetSeatAvailability}.
\end{itemize}

Confundirlos produce contratos ambiguos. Un evento no debería ser una solicitud escondida que depende de un consumidor específico para que el sistema siga correcto.

\section{Dual write y transactional outbox}

El dual write aparece cuando un servicio actualiza su base y publica un evento. Ningún orden elimina todas las ventanas de falla.

Transactional outbox guarda estado de negocio y evento pendiente en una transacción local. Un relay o CDC publica. Esto evita “base sí / evento no”, pero introduce entrega potencialmente repetida y exige deduplicación downstream.

Un evento debería incluir:
\begin{itemize}
\item event ID;
\item aggregate ID;
\item tipo y versión de esquema;
\item timestamp del dominio/producción;
\item datos mínimos necesarios;
\item metadata de tracing y causalidad cuando corresponda.
\end{itemize}

\section{Contratos y evolución}

Los eventos almacenados duran más que un despliegue. Cambiar significado rompe consumidores actuales, futuros y replay histórico.

Prácticas:
\begin{itemize}
\item cambios aditivos compatibles;
\item campos opcionales y defaults;
\item validación de schemas;
\item ownership y documentación;
\item pruebas de contrato;
\item deprecación planificada;
\item consumidores tolerantes a datos desconocidos.
\end{itemize}

\section{Sagas}

Una saga modela un workflow como transacciones locales. Si un paso falla, el sistema puede reintentar, continuar con una alternativa o compensar pasos previos.

\subsection{Compensación}

No es rollback físico. Una devolución no borra el cobro histórico; un email de cancelación no desenvía el original; liberar un cupo puede no recuperar una oportunidad perdida.

Las compensaciones deben diseñarse con semántica de negocio y ser idempotentes.

\subsection{Ejemplo: inscripción paga}

\begin{enumerate}
\item crear inscripción pendiente;
\item reservar cupo;
\item autorizar pago;
\item activar acceso;
\item emitir comprobante.
\end{enumerate}

Si falla pago, liberar cupo y cancelar. Si falla email, no cancelar acceso: reintentar un paso secundario. Esto muestra que el flujo no tiene una única política de rollback.

\section{Coreografía y orquestación}

\subsection{Coreografía}

Cada servicio reacciona a eventos. Adecuada para pocos participantes y flujos simples. El riesgo aumenta con ciclos, dependencias implícitas y dificultad de visualizar el estado global.

\subsection{Orquestación}

Un coordinador durable mantiene el estado del workflow, envía comandos y decide compensaciones. Mejora visibilidad y control, pero puede concentrar lógica y criticidad.

AWS Prescriptive Guidance documenta ambas variantes y advierte sobre dual writes, ciclos, timeouts y complejidad de compensación. El patrón no elimina estos problemas; los organiza.

\section{Strangler Fig y migración incremental}

Una migración big-bang concentra riesgo y congela entrega. Strangler Fig permite interceptar una capacidad, implementarla fuera del monolito, migrar tráfico y retirar gradualmente la versión antigua.

Una extracción debería tener hipótesis:
\begin{itemize}
\item qué problema de escalado/cambio resuelve;
\item qué SLO mejora;
\item qué blast radius reduce;
\item qué costo operativo agrega;
\item cómo se revierte si la hipótesis falla.
\end{itemize}

\section{Operabilidad como requisito previo}

Antes de multiplicar servicios, la organización necesita:
\begin{itemize}
\item CI/CD y rollback;
\item observabilidad y tracing;
\item gestión de configuración y secretos;
\item estándares de resiliencia;
\item on-call y ownership;
\item catálogo de servicios y documentación;
\item mecanismos de compatibilidad de contratos.
\end{itemize}

Sin estas capacidades, cada servicio multiplica toil y variabilidad.

\section{Matriz de decisión}

\begin{tabular}{p{0.24\textwidth}p{0.32\textwidth}p{0.34\textwidth}}
\toprule
Dimensión & Mantener junto & Separar \\
\midrule
Invariantes & transacción local frecuente & consistencia eventual/contrato estable \\
Carga & similar & escalado/costo muy distinto \\
Equipos & ownership compartido & equipos autónomos y maduros \\
Falla & debe caer como unidad & conviene aislar blast radius \\
Cambio & mismo ritmo & ciclos de despliegue distintos \\
Datos & modelo común coherente & bounded contexts distintos \\
\bottomrule
\end{tabular}

\section{Método de diseño}

\begin{enumerate}
\item modelar dominio e invariantes;
\item identificar capacidades y owners;
\item trazar flujos críticos;
\item elegir comunicación por interacción;
\item localizar dual writes y sagas;
\item diseñar fallas y degradación;
\item verificar capacidad organizacional;
\item definir métricas y criterio de extracción;
\item migrar incrementalmente;
\item revisar fronteras con evidencia.
\end{enumerate}


\section{Precisiones técnicas indispensables}

\subsection{Bounded context no implica proceso}

Un bounded context es una frontera semántica dentro de la cual el modelo tiene significado coherente. Puede implementarse como módulo de un monolito. La frontera de despliegue se justifica después, cuando autonomía, escala, seguridad o aislamiento pagan la prima distribuida.

\subsection{Database per service significa ownership, no una máquina por servicio}

El principio esencial es que un único owner modifica el modelo. La separación puede realizarse mediante esquemas, bases o clusters distintos según el riesgo. Compartir el mismo servidor físico no invalida el patrón; permitir escrituras laterales a tablas ajenas sí destruye el ownership.

\subsection{Outbox no ofrece exactly-once global}

Outbox hace atómicos el cambio local y la intención de publicar. El relay puede reenviar y el broker puede entregar nuevamente. El consumidor necesita inbox o deduplicación atómica con el efecto. Para efectos físicos o proveedores externos todavía puede ser necesaria reconciliación.

\subsection{Una saga es una máquina de estados durable}

No basta con encadenar handlers. El estado del workflow, los IDs de paso, deadlines, intentos, resultados y compensaciones deben sobrevivir reinicios. Las transiciones deben ser idempotentes y algunas fallas requieren intervención manual, no un rollback automático.

\subsection{Compensar no borra la historia}

Un reembolso es otro movimiento; un email de corrección no desenvía el anterior; devolver stock no recupera el tiempo perdido. La compensación restaura una condición de negocio aceptable, puede tardar y también puede fallar.

\subsection{La autonomía debe demostrarse}

Una extracción se justifica con evidencia: mayor frecuencia de despliegue independiente, menor lead time, menor blast radius o escalado diferenciado. Si los equipos, esquemas y releases siguen coordinándose siempre, probablemente se creó un monolito distribuido.

\section{Errores frecuentes}

\begin{itemize}
\item servicio por entidad o tabla;
\item usar microservicios como objetivo;
\item database sharing sin ownership;
\item cadenas RPC profundas;
\item event-driven sin contratos;
\item outbox sin deduplicación;
\item saga sin compensación idempotente;
\item big-bang rewrite;
\item ignorar organización y on-call;
\item medir cantidad de servicios en vez de autonomía o SLO.
\end{itemize}

\section{Bibliografía guiada y casos reales}

\begin{itemize}
\item Kleppmann, \emph{Designing Data-Intensive Applications}, capítulos 11 y 12.
\item Fowler y Lewis, \href{https://martinfowler.com/articles/microservices.html}{Microservices} y Fowler, \href{https://martinfowler.com/bliki/MonolithFirst.html}{Monolith First}.
\item Fowler, \href{https://martinfowler.com/articles/microservice-trade-offs.html}{Microservice Trade-Offs}.
\item Netflix TechBlog, \href{https://netflixtechblog.com/scaling-event-sourcing-for-netflix-downloads-episode-1-6bc1595c5595}{Scaling Event Sourcing for Netflix Downloads}.
\item AWS Prescriptive Guidance: \href{https://docs.aws.amazon.com/prescriptive-guidance/latest/cloud-design-patterns/transactional-outbox.html}{Transactional Outbox}, \href{https://docs.aws.amazon.com/prescriptive-guidance/latest/cloud-design-patterns/saga-patterns.html}{Saga Patterns} y \href{https://docs.aws.amazon.com/prescriptive-guidance/latest/cloud-design-patterns/strangler-fig.html}{Strangler Fig}.
\end{itemize}

\section{Conexión con el Trabajo Final Integrador}

Cada grupo debe justificar:
\begin{itemize}
\item qué queda unido y por qué;
\item qué capacidad se separa y qué presión lo exige;
\item ownership de datos;
\item interacciones RPC/evento;
\item dual writes y outbox;
\item saga y compensaciones;
\item capacidades operativas necesarias;
\item criterio para revisar la frontera.
\end{itemize}


\section{Ejemplos cuantitativos de la prima distribuida}

\subsection{Disponibilidad de una cadena sincrónica}

Si una request depende secuencialmente de API, A, B y C, y cada componente tiene disponibilidad independiente de 99,9\%, la disponibilidad idealizada del flujo es:
\[
A_{flujo}=0{,}999^4\approx0{,}996006=99{,}6006\%.
\]
La cuenta ignora red, timeouts y fallas correlacionadas, por lo que es optimista. Además, las latencias se suman y los percentiles de cola pueden dominar. La separación creó una cadena cuya confiabilidad debe evaluarse como flujo, no como componentes aislados.

\subsection{Costo de coordinación entre equipos}

Suponga que un cambio en Inventario requiere aprobación de tres equipos y cuatro repositorios. Aunque existan servicios separados, la unidad real de cambio sigue siendo conjunta. Métricas como lead time, porcentaje de despliegues independientes y cantidad de coordinaciones cruzadas revelan si la frontera compra autonomía o solo red.

\section{Caso real: modularidad en Shopify}

Shopify ha documentado la evolución de un gran monolito modular y herramientas para hacer explícitas dependencias entre componentes. La lección no es copiar una herramienta específica, sino observar que una única unidad de despliegue puede conservar fronteras mediante ownership, APIs internas, chequeos automatizados y reducción de dependencias cíclicas.

Una organización puede elegir un monolito modular cuando:
\begin{itemize}
\item las transacciones locales son frecuentes;
\item el dominio aún cambia;
\item el equipo no tiene capacidad operacional para múltiples servicios;
\item la separación no mejora escala, ownership o cadencia de despliegue.
\end{itemize}

\section{Dual write trabajado paso a paso}

El servicio de inscripción actualiza \texttt{Enrollment} y publica \texttt{StudentEnrolled}.

\begin{center}
\begin{tabular}{p{0.22\textwidth}p{0.31\textwidth}p{0.31\textwidth}}
\toprule
Orden & Ventana de falla & Inconsistencia \\
\midrule
DB luego broker & evento publicado, DB aborta & consumidor ve inscripción inexistente \\
DB antes broker & DB confirma, proceso cae & inscripción sin evento \\
ambos terminan, ACK se pierde & retry completo & evento o efecto duplicado \\
\bottomrule
\end{tabular}
\end{center}

Transactional outbox convierte el primer paso en una sola transacción local: negocio y evento pendiente. El relay puede repetir, por lo que cada evento necesita identidad y el consumidor debe ser idempotente.

\subsection{Schema orientativo}

\begin{verbatim}
outbox(
  event_id UUID PRIMARY KEY,
  aggregate_id TEXT NOT NULL,
  event_type TEXT NOT NULL,
  schema_version INT NOT NULL,
  payload JSONB NOT NULL,
  sequence_no BIGINT NOT NULL,
  created_at TIMESTAMP NOT NULL,
  published_at TIMESTAMP NULL
)
\end{verbatim}

Debe monitorearse edad máxima, backlog, errores del relay y duplicados detectados.

\section{Saga como máquina de estados durable}

Una saga de compra puede modelarse con estados:
\[
NEW\rightarrow STOCK\_RESERVED\rightarrow PAYMENT\_AUTHORIZED\rightarrow COMPLETED.
\]
Ante falla después del pago:
\[
PAYMENT\_AUTHORIZED\rightarrow REFUND\_PENDING\rightarrow COMPENSATED.
\]
Cada transición debe tener:
\begin{itemize}
\item identificador de saga y de paso;
\item deadline y política de retry;
\item comando idempotente;
\item resultado durable;
\item compensación o intervención manual;
\item trazabilidad.
\end{itemize}

No todas las fallas exigen volver atrás. Si solo falla el email, se reintenta la notificación sin cancelar el negocio ya confirmado.

\section{Coreografía y orquestación con un criterio de decisión}

\begin{center}
\begin{tabular}{p{0.20\textwidth}p{0.32\textwidth}p{0.34\textwidth}}
\toprule
Criterio & Coreografía & Orquestación \\
\midrule
Participantes & pocos y estables & muchos o workflow largo \\
Visibilidad & reconstruida por eventos & estado explícito de proceso \\
Cambio & local, pero contrato sensible & centraliza evolución del flujo \\
Falla & ciclos y eventos perdidos & orquestador crítico \\
Intervención & más difícil & natural con estados y comandos \\
\bottomrule
\end{tabular}
\end{center}

La elección depende de auditabilidad, duración, compensaciones, necesidad de intervención y frecuencia de cambio.

\section{Migración incremental y evidencia}

Strangler Fig usa un router/gateway para trasladar rutas gradualmente. Una extracción seria declara:

\begin{enumerate}
\item hipótesis: qué problema resolverá;
\item baseline: latencia, lead time, costo e incidentes actuales;
\item ruta inicial de bajo riesgo;
\item coexistencia y sincronización de datos;
\item rollback;
\item criterio de éxito y condición para detener la migración.
\end{enumerate}

Un big-bang reemplaza demasiadas variables a la vez y dificulta aprender qué decisión produjo el resultado.

\section{Métricas para demostrar que una frontera vale}

\begin{itemize}
\item lead time y frecuencia de despliegue por capacidad;
\item porcentaje de despliegues realmente independientes;
\item blast radius y tiempo de recuperación de incidentes;
\item p95/p99 y disponibilidad del flujo extremo a extremo;
\item cantidad de coordinaciones entre equipos por cambio;
\item backlog y edad de outbox;
\item sagas abiertas, compensadas y enviadas a intervención;
\item cambios incompatibles de contrato;
\item costo de infraestructura y toil por servicio.
\end{itemize}

\section{Glosario operativo}

\begin{description}
\item[Capacidad.] Responsabilidad de negocio con lenguaje e invariantes propios.
\item[Bounded context.] Límite dentro del cual un modelo tiene significado consistente.
\item[Ownership.] Autoridad exclusiva para modificar datos y operar una capacidad.
\item[Monolito distribuido.] Servicios separados que conservan acoplamiento de despliegue, datos y cambio.
\item[Outbox.] Registro transaccional de eventos pendientes junto con el cambio de negocio.
\item[Saga.] Workflow durable compuesto por transacciones locales y acciones de continuación/compensación.
\end{description}

\section{Autoevaluación ampliada}

\begin{enumerate}
\item ¿Qué distingue un monolito modular de una bola de barro?
\item ¿Qué significa que un servicio compre autonomía?
\item ¿Qué costos forman la prima distribuida?
\item ¿Qué presión concreta justificaría extraer Inventario?
\item ¿Por qué un servicio por tabla suele ser una mala frontera?
\item ¿Qué es un bounded context?
\item ¿Qué implica ownership de datos?
\item ¿Cómo se resuelve una consulta global sin compartir tablas?
\item ¿Qué revela Conway sobre fronteras técnicas?
\item ¿Cuándo conviene RPC?
\item ¿Cuándo conviene publicar un evento?
\item ¿Qué diferencia un comando, evento y consulta?
\item ¿Por qué ninguna orden resuelve dual write?
\item ¿Qué resuelve outbox y qué deja abierto?
\item ¿Cómo evoluciona un contrato de evento de forma compatible?
\item ¿Por qué una compensación no es rollback?
\item ¿Cuándo elegir coreografía?
\item ¿Cuándo elegir orquestación?
\item ¿Qué reduce el riesgo de una migración Strangler?
\item ¿Qué métrica demostraría autonomía real después de extraer un servicio?
\end{enumerate}

\section{Síntesis final}

\begin{uakeybox}
La arquitectura correcta no es la más distribuida. Es la que alinea fronteras de dominio, ownership, datos, despliegue y falla. Cada servicio debe comprar una autonomía real que justifique la prima distribuida que introduce.
\end{uakeybox}

\end{document}
